\documentclass[
	spanish,
	9pt,
	xcolor=table,
	handout,
	aspectratio=1610,
  ignorenonframetext
]{beamer}

\usepackage[spanish]{babel}
\usepackage[
  cursointroductorio,
  fullpagenumbering
]{dunestyle-beamer}
\usepackage{booktabs}
\usepackage{longtable}

\usepackage[backend=biber,style=numeric, defernumbers=true, sorting=ynt,maxbibnames=4,maxcitenames=4]{biblatex}
\addbibresource{references.bib}

\title{
  Introducción al módulo \texttt{dune-common}
}
\subtitle{
  Conceptos básicos:
  \lstinline{bash},
  \lstinline{g++},
  \lstinline{cmake},
  \lstinline{gnuplot},
  \lstinline{gmsh},
  \lstinline{paraview}
}

\author{}
\institute[]
{
  \noindent
  Practique los ejemplos en Gitpod

	\href{https://gitpod.io\#/https://github.com/cpp-review-dune/dune-basics}{\includegraphics[width=4.5cm]{open-in-gitpod}}	
	
  \vfill

  \begin{minipage}{0.12\paperwidth}
    Disponible en
  \end{minipage}%
  \hfill%
  \begin{minipage}{0.18\paperwidth}
    \href{https://github.com/cpp-review-dune}{\includegraphics[width=1.2cm]{GitHub-Mark-120px-plus}}
  \end{minipage}

  \vfill

  \begin{minipage}{0.26\paperwidth}
    ¡Únete al grupo en Telegram!
  \end{minipage}%
  \hfill%
  \begin{minipage}{0.04\paperwidth}
    \href{https://t.me/joinchat/OsfYP1xnFlxjN2Ix}{\includegraphics[width=1.2cm]{telegram-logo}}
  \end{minipage}
  
}